% =========================
% chapters/02_quantum_background.tex
% =========================
\chapter{Quantum Control Background for Non-Specialists}
\label{ch:quantum_background}

\section{Minimal quantum computing vocabulary}
A quantum bit (qubit) is a two-level system whose state can be written as
$\lvert \psi \rangle = \alpha \lvert 0 \rangle + \beta \lvert 1 \rangle$, with complex amplitudes satisfying
$|\alpha|^2 + |\beta|^2 = 1$. In the circuit model, computation proceeds via unitary gates and measurements.
Universal quantum computation can be achieved with a suitable set of single-qubit gates and an entangling
two-qubit gate (e.g., CNOT plus arbitrary single-qubit rotations).

% TODO: add one citation to a clean excerpt or textbook chapter you’ll cite in BibTeX:
% \cite{nielsen_chuang_ch4_excerpt}

\section{Trapped-ion qubits and gate primitives}
In trapped-ion platforms, qubits are encoded in internal electronic states of ions confined by electromagnetic
fields. Laser-mediated interactions couple internal states to collective motional modes, enabling entangling
operations.

A widely used native entangling gate is the M{\o}lmer--S{\o}rensen (MS) gate, which uses bichromatic driving
near motional sidebands to generate an effective spin--spin interaction while being relatively insensitive to
thermal motional occupation compared to earlier schemes that required ground-state cooling.

% Replace with the references in your BibTeX:
% \cite{sorensen_molmer_1999_prl,haeffner_2008_review,bruzewicz_2019_apr}

\begin{figure}[t]
  \centering
  \draftfigure[0.85\linewidth]{figures/trapped_ion_control_stack.pdf}
  \caption{Placeholder overview figure: trapped-ion control stack showing qubit states, laser beams, AOMs,
  and the real-time electronics layer. (Generate as a vector PDF.)}
  \label{fig:ti_control_stack}
\end{figure}

\section{Why RF control matters: AOMs and optical control pathways}
Acousto-optic modulators (AOMs) are commonly used to control laser intensity, switch beams, and impose
frequency shifts by diffracting light from an RF-driven acoustic wave in a crystal. In many laboratory
configurations, the RF drive amplitude sets optical diffraction efficiency (intensity control), while the RF
frequency determines the optical frequency shift and beam deflection angle (in deflector use cases).

Because quantum gates can require precise pulse timing, shaped envelopes, and stable phase relationships,
the RF generation chain (synthesizer/AWG, amplifiers, AOM driver) becomes a control bottleneck if timing
is non-deterministic or if phase is not repeatable.

\section{Control requirements implied by gates}
This thesis uses three control requirements as design drivers:
(i) deterministic timing of updates and triggers relative to a shared experiment timeline,
(ii) repeatable and characterizable phase relationships within and across channels,
and (iii) sufficient waveform flexibility (multitone and pulse shaping) for modern gate and calibration sequences.

% TODO: Replace with your experiment’s concrete requirements (e.g., pulse durations, detunings, number of tones).
